% Laget automatisk av elevark/lag_elevark.py fra opplegg/rekyl-vogner.html. Endringer her blir overskrevet.
\documentclass[11pt]{article}
\usepackage{elevark}
\begin{document}
\elevarkhode{Fysikk 1 · Bevegelsesmengde og støt}{Rekyl med to vogner}{https://fysikklab.no/opplegg/rekyl-vogner.html}{Opplegget på nett}

To vogner står i ro med en spent fjær mellom seg. Når fjæren løses ut, ruller de hver sin vei. Plasser vognene slik at de treffer hver sin endekloss samtidig. Hva sier strekningene om massene? Forsøket trenger ingen sensorer, bare øret og en linjal.

\section*{Hensikt}

\begin{itemize}
\item Bruke bevaring av bevegelsesmengde på en «eksplosjon» der alt starter i ro.
\item Se at forholdet mellom strekningene er det omvendte av forholdet mellom massene.
\end{itemize}

\section*{Utstyr}

\begin{itemize}
\item To vogner, der minst én har et fjærstempel som kan spennes og løses ut.
\item En rett bane med en endekloss i hver ende, lodd til vognene, vekt og linjal eller målebånd.
\end{itemize}

\section*{Teorien}

Før utløsningen er bevegelsesmengden null. Rett etterpå er den fortsatt null: $m_{1}v_{1} = m_{2}v_{2}$ (som størrelser, i hver sin retning). Hvis vognene bruker like lang tid $t$ til hver sin endekloss, er $v_{1} = d_{1}/t$ og $v_{2} = d_{2}/t$. Da blir $m_{1}d_{1} = m_{2}d_{2}$, altså $d_{1}/d_{2} = m_{2}/m_{1}$.

\section*{Framgangsmåte}

\begin{enumerate}
\item Sjekk at banen er vannrett: En vogn skal bli stående i ro hvor som helst på banen.
\item Vei vognene. Kall vognen med fjærstempelet vogn 1.
\item Spenn fjæren og sett vognene inntil hverandre et sted på banen. Løs ut fjæren ved å slå lett på utløseren ovenfra med en linjal, slik at dere ikke dytter vognene.
\item Lytt: Hører dere ett smell eller to? Flytt startstedet mot den vognen som kom fram først, og prøv igjen til dere hører ett smell.
\item Mål strekningene $d_{1}$ og $d_{2}$ fra fronten av hver vogn til endeklossen den treffer.
\item Legg lodd på vogn 2 og gjenta, slik at dere får minst fem ulike masseforhold.
\end{enumerate}

\section*{Oppgaver}

\begin{oppgaveliste}
\oppgave \fase{gjett}{Gjett}Vogn 2 har dobbelt så stor masse som vogn 1. Hvor må vognene starte for å treffe endeklossene samtidig, når banen mellom klossene er 100~cm lengre enn de to vognene til sammen?
\svar
\oppgave \fase{test}{Mål}Gjør målingene, eller trykk på «Fyll inn eksempel». Ligger punktene nær linjen med stigningstall 1?
\svar
\oppgave \fase{forklar}{Forklar}Hvorfor betyr «ett smell» at farten til hver vogn er proporsjonal med strekningen den kjører?
\svar
\oppgave \fase{test}{Regn}Vogn 1 har masse 0,50~kg og vogn 2 har masse 1,50~kg. Hvilken vogn får mest kinetisk energi, og hvor mange ganger så mye?
\svar
\oppgave \fase{forklar}{Vurder}Friksjonen bremser begge vognene. Hvorfor ødelegger ikke det metoden helt? Når blir feilen størst?
\svar
\oppgave \fase{utfordring}{Utfordring}Du står på et skateboard (du og brettet har til sammen 60~kg) og kaster en medisinball på 3,0~kg rett fram med 4,0~m/s. Hvor fort ruller du bakover? Hvorfor kjennes det likevel ikke som du får mye fart?
\svar
\end{oppgaveliste}

\end{document}
